\cleardoublepage

% Start with German abstracrt
\begin{otherlanguage}{ngerman}
\addchap{Kurzfassung}

\ac{KI} verändert autonome Systeme durch bedeutende Fortschritte in den Bereichen Sensorik, Kommunikation und Rechenleistung. Obwohl vollständig autonome Fahrzeuge (\acp{AV}) voraussichtlich immer häufiger werden, werden sie auf absehbare Zeit gemeinsam mit von Menschen gesteuerten Fahrzeugen (\acp{HDV}) am Straßenverkehr teilnehmen. Um in diesem gemischten Verkehrsumfeld sicher zu operieren, müssen autonome Systeme in der Lage sein, menschliche Fahrhandlungen zu modellieren und vorherzusagen.
\par Diese Arbeit adressiert dieses wichtige vorgelagerte Problem durch die Entwicklung eines transparenten, probabilistischen Rahmens zur Vorhersage hochrangiger menschlicher Fahrhandlungen unter Verwendung dynamischer Bayes’scher Netze (\acp{DBN}). Ziel dieser Forschung ist die Entwicklung eines Modells, das zukünftige Fahreraktionen abschätzen kann und gleichzeitig Unsicherheiten in einer interpretierbaren Form darstellt. Ein solches Modell kann innerhalb eines \ac{AV}-Systems als Belief-Modell dienen und Informationen für nachgelagerte Planungs- und sicherheitsrelevante Module bereitstellen.
\par Die Studie verwendet den naturalistischen Autobahnfahrdatensatz highD, um zeitliche Muster menschlichen Fahrverhaltens zu erlernen. Fahrzeugtrajektorien werden in sequenzielle, fensterbasierte Beobachtungen umgewandelt, die relevante Bewegungs- und Interaktionsinformationen über die Zeit erfassen. Diese Beobachtungen fassen zusammen, wie sich ein Fahrzeug und der umgebende Verkehr über kurze Zeitintervalle entwickeln. Anschließend wird ein \ac{DBN} trainiert, um die zeitliche Entwicklung des Fahrerverhaltens zu modellieren und Aktionswahrscheinlichkeiten aus eingehenden Beobachtungen abzuleiten. Das vorgeschlagene \ac{DBN} verwendet eine hierarchische latente Struktur, in der Aktionen von Fahrregimen abhängig sind. Genauer gesagt repräsentiert das Modell zwei latente Regime, ein Regime mit geringer Interaktion und ein Regime mit hoher Interaktion, und sagt Aktionszustände innerhalb jedes Regimes vorher. Dadurch kann das Modell sowohl das vorhergesagte Fahrmanöver als auch den Interaktionskontext darstellen, in dem es auftritt.
\par Eine zentrale Stärke dieses Ansatzes liegt in seiner transparenten probabilistischen White-Box-Struktur. Im Gegensatz zu vielen Black-Box-Lernmethoden erzeugt das \ac{DBN} seine Vorhersagen über explizite probabilistische Beziehungen, die mithilfe der Wahrscheinlichkeitstheorie interpretiert werden können. Dies ist insbesondere für sicherheitsrelevante Anwendungen von Bedeutung, bei denen Interpretierbarkeit und die Darstellung von Unsicherheiten wichtig sind. Das Modell unterstützt daher nicht nur die Vorhersageleistung, sondern auch die Erklärbarkeit, was bei der Betrachtung solcher Systeme für sicherheitskritische Verkehrsanwendungen von Vorteil ist. Die Studie leistet einen Beitrag zur Forschung im Bereich des autonomen Fahrens, indem sie Aktionsvorhersage, zeitliche probabilistische Schlussfolgerung und Interpretierbarkeit in einem einheitlichen Rahmen kombiniert.
\par Die Ergebnisse zeigen, dass das Modell in der Lage ist, aussagekräftige Verhaltensmuster aus naturalistischen Verkehrsdaten zu lernen und vielversprechende Manövervorhersagen über einen zukünftigen Zeithorizont zu generieren. In der aktuellen Evaluation erreichte das Modell eine exakte Next-Step-Genauigkeit von 46{,}4\,\%, eine Makropräzision von 44{,}3\,\% sowie eine Trefferquote von 68{,}0\,\% über einen Vorhersagehorizont von zwei Sekunden. Eine Vorhersage zum Zeitpunkt \(t\) wird dabei als Treffer gezählt, wenn das vorhergesagte Manöver innerhalb der nächsten fünf Fensterschritte auftritt, was etwa zwei Sekunden entspricht. Die durchschnittliche Time-to-Hit betrug 0{,}50\,s und der Median 0{,}40\,s, was zeigt, dass viele korrekte Vorhersagen früh innerhalb des Vorhersagehorizonts auftreten. Diese Ergebnisse deuten darauf hin, dass das Modell zukünftige menschliche Fahreraktionen innerhalb eines praktisch nutzbaren kurzfristigen Zeitintervalls antizipieren kann, anstatt nur einen einzelnen exakten zukünftigen Zeitpunkt vorherzusagen. Implementierungs-Benchmarks zeigen zudem geringe Online-Inferenzkosten mit einer mittleren CPU-Latenz von 0{,}164\,ms pro Fahrzeug für Belief-Update und nächste Aktionsvorhersage.
\par Insgesamt zeigen die Ergebnisse, dass \acp{DBN} ein vielversprechender Ansatz für die probabilistische Vorhersage menschlicher Fahreraktionen sind. Die Studie demonstriert, dass hochrangige menschliche Fahrmanöver mithilfe naturalistischer Verkehrsdaten in einer interpretierbaren und zeitlich strukturierten Weise modelliert werden können. Dies stellt einen wichtigen Schritt hin zu einem besseren Verständnis menschlichen Fahrverhaltens im gemischten Verkehr dar und kann \ac{AV}-Systeme dabei unterstützen, das Verhalten umliegender menschlicher Fahrer besser einzuschätzen.

\vfill
\noindent\textbf{Stichwörter:}  Vorhersage menschlichen Fahrverhaltens, Autonomes Fahren, Von Menschen gesteuerte Fahrzeuge, Dynamische Bayes’sche Netze, Manövervorhersage, Probabilistische Modellierung, Gemischte Verkehrsumgebungen
\vfill
\acresetall
\end{otherlanguage}
% Then continue with the english one.
\begin{otherlanguage}{english}
\addchap{Abstract}

\ac{AI} is transforming autonomous systems through major advances in sensing, communication, and computation. Although fully \acp{AV} are expected to become more common, they will share the road with \acp{HDV} for the foreseeable future. To operate safely in this mixed-traffic setting, autonomous systems must be able to model and predict human driver actions.
\par This work addresses this important upstream problem by developing a transparent, probabilistic framework for predicting high-level human driving actions using \acp{DBN}. The purpose of this research is to develop a model that can estimate future driver actions while also representing uncertainty in an interpretable form. Such a model can act as a belief model inside an \ac{AV} framework and provide information for downstream planning and safety-related modules.
\par The study uses the highD naturalistic highway driving dataset to learn temporal patterns of human driving behavior. Vehicle trajectory data are converted into sequential window-based observations that capture relevant motion and interaction information over time. These observations summarize how the vehicle and its surrounding traffic evolve over short time intervals. A \ac{DBN} is then trained to model the temporal evolution of driver behavior and to infer action probabilities from incoming observations. The proposed \ac{DBN} uses a hierarchical latent structure in which actions are conditioned on driving regimes. More specifically, the model represents two latent regimes, low-interaction and high-interaction, and predicts action states within each regime. This allows the framework to represent both the predicted maneuver and the interaction context in which it occurs. 
\par A key strength of this approach is that it provides a transparent, white-box probabilistic structure. Unlike many black-box learning methods, the \ac{DBN} makes predictions through explicit probabilistic relationships that can be interpreted using probability theory. This is valuable in safety-relevant applications, where interpretability and uncertainty representation are important. The model therefore supports not only prediction, but also explainability, which is useful when such systems are considered for use in safety-critical transport settings. The study contributes to science and engineering by combining action prediction, temporal probabilistic reasoning, and interpretability in a single framework for autonomous driving applications.
\par The results show that the model can learn meaningful behavioral patterns from naturalistic traffic data and generate promising maneuver predictions over a future horizon.  In the current evaluation, the model achieved 46.4\% exact next-step accuracy, 44.3\% macro precision, and a 68.0\% hit rate over a 2-second prediction horizon. Here, a prediction made at time t is counted as a hit if the predicted maneuver appears at any point within the next 5 window steps, corresponding to about 2 seconds. The mean time-to-hit was 0.50 s, and the median time-to-hit was 0.40 s, showing that many correct predictions occur early within the horizon. These results suggest that the model can anticipate future human driver actions within a practically useful short-term interval rather than only at one exact future instant. Implementation benchmarks also indicate low online inference cost, with a mean single-vehicle \ac{CPU} latency of 0.164 ms for belief update and next-action prediction.
\par Overall, the findings indicate that \acp{DBN} are a promising framework for probabilistic human driver action prediction. The study shows that high-level human driving maneuvers can be modeled in an interpretable and temporally structured way using naturalistic traffic data. This provides a useful step toward better behavior understanding in mixed traffic and may support \ac{AV} systems that need to reason about nearby human drivers.


\vfill
\noindent\textbf{Keywords:} Human Driver Behavior Prediction, Autonomous Driving, Human-Driven Vehicles, Dynamic Bayesian Networks, Maneuver Prediction, Probabilistic Modeling, Mixed Traffic Environments
\vfill
\end{otherlanguage}